
\subsubsection{6.1 Key Findings Interpretation}

Validation demonstrates that the infrastructure layer (measurement and mechanisms) works at proof-of-concept scale. All five sub-hypotheses passed gates, confirming that automated health metrics can detect deprecation candidates (88.3\% precision, 100\% recall), context-aware graphs can provide task-specific recommendations (100\% accuracy on synthetic patterns), and load-time instrumentation can track adoption events with minimal overhead (0.21-5\%). Three-component synergy was validated: health metrics surface candidates automatically, context graphs personalize recommendations, and instrumentation enables outcome measurement.

Results exceed targets across all metrics (precision 88.3\% versus 60\% target, overhead 0.21\% versus 10\% target), but perfect scores (100\% recall, 100\% context accuracy) are likely optimistic due to synthetic and simulated data. Real-world accuracy is expected to regress to 70-95\% based on pattern cleanliness assumptions. Proof-of-concept scale (100-1000 samples) validates mechanistic feasibility but not production-scale performance (60,000+ datasets).

\subsubsection{6.2 Limitations}

\textbf{Proof-of-Concept Scale versus Production:} Experiments validated mechanisms on 100-1000 samples; production deployment (60,000 HuggingFace datasets) may reveal scalability issues. NetworkX graph queries scale to thousands of nodes, but 60,000-node performance is untested. This is an acceptable trade-off: mechanistic validation scopes to feasibility demonstration, not deployment readiness.

\textbf{Mock and Simulated Data versus Real API:} h-m2, h-m3, and h-m4 used mock data; h-m1 used simulated HuggingFace metadata. This isolates mechanism validation from external dependencies (API rate limits, authentication) but leaves generalization to real API data unvalidated. Patterns were derived from software package manager literature to ensure realism, but real-world edge cases (e.g., dataset cards without citations, ambiguous task contexts) are untested.

\textbf{Efficacy Untested:} Phase 5 baseline comparison was skipped; the P2 prediction ($\geq 50$\% adoption lift versus informal mechanisms) remains untested. Results validate the mechanistic layer (does it work?) but not the efficacy layer (how well?). This requires a controlled experiment with treatment and control groups to measure adoption lift. This is acceptable for infrastructure research: demonstration precedes deployment-scale efficacy measurement.

\textbf{Synthetic Patterns:} h-m2 context inference was validated on clean synthetic patterns (exact-match task labels, well-formed dataset cards). Real-world usage may include ambiguous contexts (multi-task datasets, missing task annotations), causing accuracy to regress below 100\%. The override mechanism (27.95\% rate) provides a safety valve but does not eliminate misclassification risk.

\subsubsection{6.3 Broader Impact}

\textbf{Positive:} Formal deprecation mechanisms improve dataset quality signals for practitioners, reduce maintainer burden (automated candidate detection versus manual triage), and enable reproducibility (users are notified of deprecated datasets at load time). Telemetry infrastructure enables longitudinal studies of deprecation efficacy (A/B testing interventions).

\textbf{Negative:} Load-time telemetry raises privacy concerns (mitigated via SHA256 hashing with no user personally identifiable information, opt-in consent). Automated deprecation warnings could disrupt workflows if false-positive rate is high (mitigated via 88.3\% precision, user override mechanism). Successor recommendations may introduce bias if context inference systematically fails for underrepresented tasks (requires fairness audits on production data).

\textbf{Equitable:} Infrastructure benefits all HuggingFace users equally; there is no differential impact by demographic. Deployment requires equitable access to instrumented loaders (maintained as open-source, not platform-gated).

\subsubsection{6.4 Generalizability}

Design patterns (health metrics, context graphs, load-time instrumentation) are applicable beyond HuggingFace. OpenML, Kaggle, Papers with Code, and UCI Machine Learning Repository share common infrastructure patterns (programmatic loaders, metadata APIs, version tracking). Context inference patterns (dataset card citation analysis, import history introspection) transfer to platforms with structured metadata. Cross-platform validation is future work.

Health metric weights (velocity=2.0, emergence=1.5, issue\_ratio=1.0) were tuned for HuggingFace usage distributions and require recalibration for platforms with different download patterns (e.g., Kaggle competitions versus academic datasets). The threshold ($\geq 2.5)$ provides a precision-recall tradeoff; production deployment benefits from platform-specific tuning.

\subsubsection{6.5 Theoretical Contribution}

This work reframes datasets as dependencies with lifecycle management needs distinct from both static files (require deprecation mechanisms) and software packages (require task-aware succession, usage-based health signals). The three-failure-mode decomposition (detection, recommendation, measurement) organizes infrastructure requirements. Sub-hypothesis decomposition methodology is applicable to other machine learning infrastructure validation (model registries, experiment tracking, deployment pipelines).
